\originalsection{Galilei变换、质心系与能量函数}
自由粒子的拉格朗日量为何取速度二次式, 还可从原8.223第4讲的对称要求看出. 空间与时间均匀性使$L$只依赖速度, 各向同性使$L=f(v^2)$. 再要求任意匀速参考系变换只使$L$改变全时间导数, 则$L(\vect v+\vect V)-L(\vect v)$至多为速度的线性式加常数. 对速度求两次导数, 有$D_v^2L(\vect v+\vect V)=D_v^2L(\vect v)$, 任意$\vect V$使Hessian为常矩阵. 各向同性又使该矩阵为$mI$, 于是$L=m v^2/2+$常数, 可能的固定线性速度项被各向同性排除或作为端点项删去. 取$m>0$是通常正动能模型的物理要求. 这是非相对论惯性系和端点不变性的推导, 不能直接用于相对论自由粒子.

任意多质点系统有$\vect r_a=\vect R+\vect\rho_a$, $\sum_am_a\vect\rho_a=0$, 从而$T=M\dot{\vect R}^2/2+\sum_am_a\dot{\vect\rho}_a^2/2$. 各项质量为正, 所以在与系统作匀速相对平移的惯性系中, 质心系给最小动能; 这里比较的是同一瞬时物理状态在不同参考系的动能. 内部运动动能不因共同匀速平移而改变.

一维匀速变换$x=X+v_0t$把自由粒子拉格朗日量变成$m\dot X^2/2+mv_0\dot X+mv_0^2/2$. 后两项为全导数$\dd(mv_0X+mv_0^2t/2)/\dd t$, 故运动方程仍为$\ddot X=0$. 对只依赖相对位置的多体势同样成立. 全导数可以改变共轭动量及能量函数的表达, 不改变相同坐标变量的运动方程.

为了看清能量输入, 令弹簧端点以$v_0$匀速运动, $L=m\dot x^2/2-k(x-v_0t)^2/2$. 惯性系机械能$E=m\dot x^2/2+k(x-v_0t)^2/2$满足$\dot E=-kv_0(x-v_0t)$, 是外部驱动端点的输入功率. 取$X=x-v_0t$, 不删全导数时$P=m(\dot X+v_0)$, 新哈密顿量为
\[
H'=\frac{(P-mv_0)^2}{2m}+\frac{kX^2}2-\frac{mv_0^2}2
=E-v_0P.
\]
它不显含时间而守恒, 与惯性系$E$的变化并不矛盾. 若再删全导数, 共轭动量改成$m\dot X$, 新能量函数改成$m\dot X^2/2+kX^2/2$. 每一步都须同时说明坐标和拉格朗日量的选择.

\originalsection{齐次势的尺度关系}
设$q$为笛卡尔位置坐标, $U(\lambda q)=\lambda^kU(q)$, 动能是固定质量的速度二次式. 若$q(t)$是解, 取$\widetilde q(t)=\lambda q(t/\lambda^{1-k/2})$. 加速度缩放为$\lambda^{k-1}$, 而$\nabla U(\lambda q)=\lambda^{k-1}\nabla U(q)$, 所以新路径仍满足牛顿方程. 周期或对应运动时长因而按$\lambda^{1-k/2}$缩放. 谐振子$k=2$周期不随振幅改变; Kepler势$k=-1$给$T^2\propto a^3$; 匀强重力$k=1$给落体时间与高度平方根成正比.

这个变换把不同大小与能量的解联系起来, 作用量一般整体乘$\lambda^{1+k/2}$, 并非总是Noether定理要求的作用量不变. 因而不能只见动力学具有尺度关系就宣布又有一个普通Noether守恒量. 位力定理也需要上一节的长时平均条件, 二者虽同用齐次势, 逻辑不同.
